\documentclass[../../../include/open-logic-section]{subfiles}

\begin{document}

\olfileid[id]{his}{bio}{chu}

\olsection{Alonzo Church}

\olphoto{church-alonzo}{Alonzo Church}

Alonzo Church lahir di Washington, DC, pada 14 Juni 1903. Ketika ia masih
kecil, sebuah insiden dengan senapan angin membuat Church buta pada satu
mata. Ia menamatkan sekolah persiapan di Connecticut pada tahun 1920 dan
memulai pendidikan universitasnya di Princeton pada tahun yang sama. Ia
menyelesaikan studi doktoralnya pada tahun 1927. Setelah beberapa tahun di
luar negeri, Church kembali ke Princeton. Church dikenal sangat sopan dan
cermat. Tulisan papan tulisnya tanpa cela, dan ia mengawetkan makalah penting
dengan melapisinya secara hati-hati menggunakan semen Duco (perekat bening).
Di luar kegiatan akademiknya, ia gemar membaca majalah fiksi ilmiah dan
tidak segan menulis kepada para editor jika menemukan ketidakakuratan dalam
tulisan mereka.

Prestasi akademik Church sangat besar. Bersama para mahasiswanya, Stephen
Kleene dan Barkley Rosser, ia mengembangkan teori keterhitungan efektif,
yaitu kalkulus lambda, secara independen dari pengembangan mesin Turing oleh
Alan Turing. Kedua definisi komputabilitas tersebut ekuivalen dan melahirkan
apa yang sekarang dikenal sebagai \emph{Tesis Church--Turing}: suatu fungsi
pada bilangan asli dapat dihitung secara efektif jika dan hanya jika fungsi
itu dapat dihitung dengan mesin Turing (atau kalkulus lambda). Ia juga
membuktikan hasil yang sekarang dikenal sebagai \emph{Teorema Church}:
masalah keputusan bagi validitas rumus orde pertama tidak dapat diselesaikan.

Church terus berkarya hingga usia lanjut. Pada tahun 1967 ia meninggalkan
Princeton menuju UCLA, tempat ia menjadi profesor sampai pensiun pada tahun
1990. Church meninggal pada 1 Agustus 1995 dalam usia 92 tahun.

\begin{reading}
Untuk biografi singkat Church, lihat \citet{EndertonND}. Tulisan asli Church
tentang kalkulus lambda dan Entscheidungsproblem (Tesis Church) adalah
\citet{Church1936,Church1936a}. \citet{Aspray1984} merekam sebuah wawancara
dengan Church mengenai komunitas matematika Princeton pada dasawarsa 1930-an.
Church menulis serangkaian ulasan buku untuk \emph{Journal of Symbolic Logic}
sejak 1936 hingga 1979. Semuanya diarsipkan di situs web John MacFarlane
\citep{MacFarlane2015}.
\end{reading}

\end{document}
